% Generated draft: requires mathematical and source review.
\documentclass[11pt]{article}
\usepackage[margin=25mm]{geometry}
\usepackage{fontspec}
\setmainfont{Latin Modern Roman}
\usepackage{amsmath,amssymb,mathtools,unicode-math}
\setmathfont{Latin Modern Math}
\usepackage{enumitem,needspace,xurl,hyperref}
\setlength{\parindent}{0pt}
\setlength{\parskip}{5pt}
\setlength{\emergencystretch}{3em}
\newcommand{\N}{\mathbb{N}}
\newcommand{\Z}{\mathbb{Z}}
\newcommand{\Q}{\mathbb{Q}}
\newcommand{\R}{\mathbb{R}}
\newcommand{\C}{\mathbb{C}}
\newcommand{\F}{\mathbb{F}}
\newcommand{\A}{\mathbb{A}}
\newcommand{\PP}{\mathbb{P}}
\newcommand{\E}{\mathbb{E}}
\newcommand{\eps}{\varepsilon}
\newcommand{\dd}{\,\mathrm{d}}
\begin{document}
\section*{UnsolvedMath ID: 20002822}
\textbf{Source problem code:} catalogue entry
\subsection{Definitions and mathematical statement}
For a prime power $q$, let $U_n(\F_q)$ be the upper unitriangular subgroup and $B_n(\F_q)$ the upper triangular invertible subgroup of $\operatorname{GL}_n(\F_q)$. The workshop asks to carry the $U_n$ supercharacter constructions over to $B_n$: is there a similar ``(module to a) Hopf algebra''? The supplied source verification treats this parenthetical as an editorial ambiguity already present in the source. Accordingly, retain both targets: a Hopf algebra of Borel superclass functions, and a module attached to an appropriate Hopf algebra. The phrase ``all of the $U_n$ work'' does not identify a finite list of constructions to extend.

The aid reports that Panov's 2018 work supplies a natural supercharacter theory of $B_n(\F_q)$ for which
\[
H=\bigoplus_{n\geq0}\operatorname{SCF}(B_n(\F_q))
\]
is a graded Hopf algebra isomorphic to partially symmetric functions in noncommuting variables. Here $\operatorname{SCF}$ means the complex span of superclass indicators for that particular theory, and the degree-zero term is $\C$. Neither the full supercharacter theory nor formulas for the Hopf operations are included in this record; arbitrary induction and restriction operations are not implicitly asserted to give this structure.

For that reported Hopf algebra, the further assertions in the aid are as follows. Its rigged-partition basis has indices describable by a subset $R\subseteq[n]$ of nonidentity torus decorations, a choice of label in $\F_q^\times\setminus\{1\}$ at each point of $R$, and a set partition of $[n]\setminus R$. Let $H_{n,r}$ be the span of basis elements with $|R|=r$. The claimed bigrading satisfies
\[
H=\bigoplus_{n\geq0}\bigoplus_{r=0}^n H_{n,r},\qquad
H_{n,r}H_{m,s}\subseteq H_{n+m,r+s},\qquad
\Delta H_{n,r}\subseteq\bigoplus_{a+b=n,\,u+v=r}H_{a,u}\otimes H_{b,v}.
\]
Writing $\operatorname{Bell}(k)$ for the number of set partitions of a $k$-element set, including $\operatorname{Bell}(0)=1$, the reported dimensions are
\[
\dim_\C H_{n,r}=\binom nr(q-2)^r\operatorname{Bell}(n-r).
\]
For $q=2$ the only nonzero torus degree is $r=0$; the factor with exponent zero is interpreted as $1$.

Let $A=\operatorname{NCSym}$ be the Hopf algebra of symmetric functions in noncommuting variables, identified in the aid with the diagonal-averaged unitriangular Hopf algebra. The aid asserts a Hopf inclusion $i:A\longrightarrow H$ onto $\bigoplus_n H_{n,0}$ and a Hopf projection $p:H\longrightarrow A$ satisfying
\[
p\circ i=\operatorname{id}_A,\qquad p(H_{n,r})=0\quad(r>0).
\]
In its function model, $p$ evaluates all nonidentity torus letters at zero. Thus $H$ has the left $A$-module structure $a\cdot h=i(a)h$ and left $A$-comodule structure $(p\otimes\operatorname{id})\Delta$. The corresponding asserted Radford biproduct is
\[
H\simeq K\mathbin{\#} A,\qquad
K=\{h\in H:(\operatorname{id}\otimes p)\Delta(h)=h\otimes1\},
\]
where $K$ carries the associated braided Hopf structure. These are supplied literature and research-aid assertions, requiring verification of the precise basis, operations and identifications.
\subsection{Short English statement}
Extend unitriangular supercharacter constructions to finite Borel groups. The source's ambiguous wording asks for either a Hopf algebra or an associated module; the supplied aid reports a Borel Hopf algebra and a split projection to $\operatorname{NCSym}$, with a refined dimension formula.
\subsection{Sources}
\begin{itemize}
\item Local UnsolvedMath registry, permanent ID $20002822$, source code \texttt{catalogue entry}; canonical \texttt{aim-representation-theory-notes.json}, \texttt{notes[93]}.
\item AIM Problem Lists, Representation theory, workshop \emph{Supercharacters and combinatorial Hopf algebras}, Problem 1.10, supplied PDF source: \url{https://aimath.org/WWN/supercharacters/supercharacters.pdf}. No remarks or original literature references are supplied. The source-verification field reports that ``(module to a)'' is original wording; the registry statement status is \texttt{reconstructed\_unverified}.
\item Panov's 2018 work is cited by name and year in the supplied research summary for the Borel Hopf algebra and partially symmetric function identification; no title or complete bibliographic citation is supplied.
\item Research attempt 1 is labeled \texttt{partial\_result}, problem status \texttt{solved}, classification \texttt{SOLVED-IN-LITERATURE}, rigor \texttt{proved}, and novelty confidence low. Its bigrading, projection and biproduct claims are machine-generated research aid, not peer reviewed, and require independent verification before citation.
\end{itemize}
\end{document}
